
\chapter{内容应用与管理平台测试}

本章按“测试环境—单元测试—功能测试—非功能与性能测试”的结构展开，对内容应用与管理平台进行测试说明。测试目标包括：验证平台核心功能满足第三章需求分析中定义的业务用例；验证关键技术链路（发布口径、权限过滤、事件驱动索引、RAG问答与治理闭环）在典型场景下能够稳定运行；并从一致性、安全与权限、可扩展性以及性能与稳定性等维度对系统的非功能需求进行验证。

\section{系统测试环境}

平台在实现上由内容管理、内容应用、内容治理与数据洞察等模块构成，依赖关系型数据库与文档数据库存储内容资产，并在内容发布后构建全文索引与向量索引以支撑检索与智能问答。由于论文工程无法直接提供企业内部真实部署细节，本节从“部署形态+关键依赖组件+测试工具链”三个维度描述测试环境口径，便于复现测试结论。平台测试环境描述如表\ref{tab:test_env}所示。

\begin{table}[!htbp]
  \centering
  \caption{测试环境描述表}
  \label{tab:test_env}
  \small
  \begin{tabularx}{\linewidth}{c|>{\raggedright\arraybackslash}X}
    \hline
    项目 & 说明 \\ \hline
    部署形态 & 服务以容器化方式部署于Linux环境，按模块拆分并独立扩缩容。 \\ \hline
    开发环境 & 开发与调试环境为macOS或Linux，使用浏览器与接口调试工具完成联调。 \\ \hline
    数据存储 & 关系型数据库用于内容元数据、权限与任务数据；文档数据库用于案例详情等灵活结构字段。 \\ \hline
    检索与向量索引 & 全文检索索引用于关键词召回与高亮展示；向量索引用于语义召回与长尾覆盖，二者共同支撑RAG。 \\ \hline
    消息通知与异步链路 & 内容发布后通过消息队列触发索引增量更新与数据回收等异步任务。 \\ \hline
    外部能力 & 智能问答与治理判别依赖大模型服务，图文与图片类检索依赖多模态embedding能力。 \\ \hline
    测试工具 & 前端功能测试使用Chrome；接口与链路测试使用Postman；并发与时延压测使用JMeter；必要时结合日志与trace\_id进行链路定位。 \\ \hline
  \end{tabularx}
\end{table}

\section{单元测试}

单元测试用于在最小可控范围内验证模块内部逻辑正确性，并通过Mock隔离外部依赖以降低测试不稳定因素。平台的单元测试重点覆盖三类逻辑：第一类是权限裁决与状态机约束（例如RBAC与content\_auth匹配、审批状态流转）；第二类是内容处理与索引构建的可重复性（例如切片构建规则、slice\_id稳定生成）；第三类是治理与洞察的关键算法单元（例如TopK候选生成、任务幂等upsert与复检回收）。

单元测试的基本原则如下：
（1）边界清晰。对外部服务（数据库、索引服务、消息队列与大模型服务）通过Mock替身模拟输入输出，仅验证当前模块的业务逻辑与异常处理。
（2）可重复执行。对时间、随机数与外部依赖进行固定化处理，避免测试随运行环境变化而不稳定。
（3）覆盖关键分支。对成功路径、非法状态、权限拒绝、幂等重试与异常退化等关键分支分别构造用例，降低线上故障风险。

\section{功能测试}

功能测试用于验证平台对外可见能力与关键链路在真实交互下满足预期。本节基于第三章的核心用例集合，对内容管理、内容应用、内容治理与数据洞察模块设计功能测试用例。测试用例采用“步骤-预期结果”形式描述，覆盖正常流程、非法状态、权限约束、异步一致性与退化策略等场景。

\subsection{内容管理功能测试}

创建空间与应用是平台组织与权限的起点，用例用于验证空间隔离与应用边界能够正确生效。创建空间/应用的测试用例如表\ref{tab:tc01}所示。

\begin{table}[!htbp]
  \centering
  \caption{创建空间/应用测试用例}
  \label{tab:tc01}
  \small
  \begin{tabularx}{\linewidth}{c|>{\raggedright\arraybackslash}X}
    \hline
    用例编号 & TC01 \\ \hline
    测试内容 & 创建空间与应用 \\ \hline
    测试步骤 & 1) 进入空间管理创建空间并填写负责人。\newline 2) 在空间内创建应用并选择支持的内容形态与组织方式。\newline 3) 进入成员与角色页面绑定默认角色。 \\ \hline
    预期结果 & 1) 空间与应用创建成功并可查询详情。\newline 2) 默认角色初始化完成且创建人具备对应权限。\newline 3) 不同空间的数据互相隔离。 \\ \hline
  \end{tabularx}
\end{table}

内容生产链路覆盖创建、编辑、提审、审批与发布等过程，且要求读链路对外仅展示已发布内容。内容发布审批的测试用例如表\ref{tab:tc02}所示。

\begin{table}[!htbp]
  \centering
  \caption{内容提审与审批发布测试用例}
  \label{tab:tc02}
  \small
  \begin{tabularx}{\linewidth}{c|>{\raggedright\arraybackslash}X}
    \hline
    用例编号 & TC02 \\ \hline
    测试内容 & 内容提审、审批通过/驳回与发布口径 \\ \hline
    测试步骤 & 1) 创建一条knowledge内容并保存草稿。\newline 2) 发起提审进入待审批状态。\newline 3) 审批人执行驳回并填写意见。\newline 4) 再次提审并审批通过，内容进入已发布状态。\newline 5) 在门户与检索侧分别访问该内容。 \\ \hline
    预期结果 & 1) 草稿状态不对外可见。\newline 2) 驳回后内容回到可编辑状态并保留意见。\newline 3) 审批通过后内容对外可见，门户与检索均仅返回已发布口径。 \\ \hline
  \end{tabularx}
\end{table}

\subsection{内容应用功能测试}

内容检索用于验证多字段召回、高亮与过滤逻辑是否正确。内容检索测试用例如表\ref{tab:tc04}所示。

\begin{table}[!htbp]
  \centering
  \caption{内容检索测试用例}
  \label{tab:tc04}
  \small
  \begin{tabularx}{\linewidth}{c|>{\raggedright\arraybackslash}X}
    \hline
    用例编号 & TC04 \\ \hline
    测试内容 & 内容检索与高亮展示 \\ \hline
    测试步骤 & 1) 在应用内准备多条不同类型内容并发布。\newline 2) 使用关键词分别命中title、summary与正文切片。\newline 3) 设置类目/标签/类型过滤并分页查询。\newline 4) 对不同权限用户重复检索。 \\ \hline
    预期结果 & 1) 关键词命中不同字段均可召回且排序合理。\newline 2) 返回结果包含高亮片段与摘要。\newline 3) 结构化过滤与分页生效。\newline 4) 越权内容不会出现在检索结果中。 \\ \hline
  \end{tabularx}
\end{table}

智能问答用于验证“问题改写+多通道召回+重排+引用输出”的端到端链路，以及未命中与超时场景下的退化策略。智能问答测试用例如表\ref{tab:tc05}所示。

\begin{table}[!htbp]
  \centering
  \caption{智能问答测试用例}
  \label{tab:tc05}
  \small
  \begin{tabularx}{\linewidth}{c|>{\raggedright\arraybackslash}X}
    \hline
    用例编号 & TC05 \\ \hline
    测试内容 & RAG智能问答端到端链路 \\ \hline
    测试步骤 & 1) 输入口语化问题触发问题改写。\newline 2) 观察ES与向量两路召回候选数量与MinScore过滤。\newline 3) 检查重排后TopN证据集合与content聚合情况。\newline 4) 生成回答并检查引用映射到content与slice。\newline 5) 构造未命中与大模型超时场景。 \\ \hline
    预期结果 & 1) 改写后召回稳定性提升且不改变问题语义。\newline 2) 证据集合与问题强相关并包含可追溯引用。\newline 3) 未命中返回补充建议并进入覆盖度候选。\newline 4) 超时退化为检索摘要返回，且权限过滤始终生效。 \\ \hline
  \end{tabularx}
\end{table}

\subsection{内容治理与数据洞察功能测试}

重复/歧义治理用于验证候选召回、LLM判别与任务生成的闭环过程。重复/歧义治理测试用例如表\ref{tab:tc06}所示。

\begin{table}[!htbp]
  \centering
  \caption{重复/歧义治理测试用例}
  \label{tab:tc06}
  \small
  \begin{tabularx}{\linewidth}{c|>{\raggedright\arraybackslash}X}
    \hline
    用例编号 & TC06 \\ \hline
    测试内容 & 重复/歧义检测与子任务生成 \\ \hline
    测试步骤 & 1) 在空间内准备两条高度相似内容并发布。\newline 2) 创建治理主任务并触发重复/歧义检测计算。\newline 3) 查看生成的重复/歧义子任务与内容对。\newline 4) 修改其中一条内容并重新发布，等待日巡检复检。 \\ \hline
    预期结果 & 1) 候选召回能够命中相似内容对并生成子任务。\newline 2) 子任务幂等写入，不重复生成。\newline 3) 内容修订后复检能够关闭已消除问题的子任务。 \\ \hline
  \end{tabularx}
\end{table}

覆盖度治理用于验证“洞察聚类筛选薄弱主题$\rightarrow$生成覆盖任务$\rightarrow$补充内容$\rightarrow$复检回收”的闭环。覆盖度治理测试用例如表\ref{tab:tc07}所示。

\begin{table}[!htbp]
  \centering
  \caption{覆盖度治理测试用例}
  \label{tab:tc07}
  \small
  \begin{tabularx}{\linewidth}{c|>{\raggedright\arraybackslash}X}
    \hline
    用例编号 & TC07 \\ \hline
    测试内容 & 覆盖度任务生成与复检回收 \\ \hline
    测试步骤 & 1) 在问答链路制造未命中问题并沉淀为会话数据。\newline 2) 在洞察侧聚类生成低解决率主题。\newline 3) 在治理侧创建主任务并生成覆盖子任务。\newline 4) 补充相关内容并发布。\newline 5) 等待巡检复检回收。 \\ \hline
    预期结果 & 1) 覆盖任务能够以cluster\_id或代表性query定位问题主题。\newline 2) 内容补充发布后复检能够将覆盖任务自动关闭。\newline 3) 任务状态与处理人记录完整可追溯。 \\ \hline
  \end{tabularx}
\end{table}

数据洞察用于验证事件规范、trace\_id链路追踪与下钻能力。数据洞察测试用例如表\ref{tab:tc08}所示。

\begin{table}[!htbp]
  \centering
  \caption{数据洞察链路测试用例}
  \label{tab:tc08}
  \small
  \begin{tabularx}{\linewidth}{c|>{\raggedright\arraybackslash}X}
    \hline
    用例编号 & TC08 \\ \hline
    测试内容 & trace\_id链路追踪与会话下钻 \\ \hline
    测试步骤 & 1) 发起一次问答请求并记录trace\_id。\newline 2) 校验rag\_recall与llm\_generate等事件明细携带相同trace\_id与session\_id。\newline 3) 在洞察侧按聚类主题下钻到会话明细。\newline 4) 按trace\_id还原证据集合与引用内容。 \\ \hline
    预期结果 & 1) 关键事件均被采集且字段口径一致。\newline 2) 可按trace\_id复现单次问答的召回与生成链路。\newline 3) 聚类下钻可定位高频主题与薄弱点，为治理任务生成提供输入。 \\ \hline
  \end{tabularx}
\end{table}

\section{非功能与性能测试}

第三章从一致性与可追溯性、安全与权限、可扩展性、性能与稳定性四个方面给出了系统非功能需求。本节在前述单元测试与功能测试基础上，进一步补充对应的验证口径与性能结果，其中性能验证重点聚焦于平台实际高频链路，包括内容列表、内容检索、内容详情、智能问答，以及内容发布后的索引构建链路。

\subsection{性能测试参数说明}

为保证性能测试结果具有可解释性，本文在模拟环境中固定测试数据集、压测工具与采样口径。测试数据集包含10万级已发布内容，覆盖文章、问答、课程和案例四类内容形态。接口压测采用JMeter生成并发请求，单轮压测持续10分钟，前2分钟作为预热阶段，后8分钟计入统计结果。内容列表、内容检索和内容详情接口采用50个并发线程，智能问答接口采用20个并发线程；每类接口重复执行3轮，并取平均结果作为表中数据。测试服务以容器化方式部署，接口服务配置为8核16GB内存，检索与向量索引组件分别独立部署，压测过程中记录平均时延、P95、P99、错误率、消息积压和资源使用情况。

\subsection{非功能需求验证}

非功能需求验证结果如表\ref{tab:nfr_verify}所示。整体上，平台在一致性、安全隔离、扩展演进以及高频读写链路稳定性方面均满足论文设计目标。

\begin{table}[!htbp]
  \centering
  \caption{非功能需求验证结果表}
  \label{tab:nfr_verify}
  \small
  \begin{tabularx}{\linewidth}{>{\raggedright\arraybackslash}p{0.19\linewidth}|>{\raggedright\arraybackslash}p{0.36\linewidth}|>{\raggedright\arraybackslash}X}
    \hline
    非功能需求 & 验证方式 & 结果摘要 \\ \hline
    一致性与可追溯性 & 结合内容提审、审批、发布、索引更新与问答事件，验证发布口径、事件消费与trace\_id链路还原。 & 门户、检索与问答链路均仅使用已发布内容；可按trace\_id还原单次问答的召回、重排与生成过程。 \\ \hline
    安全与权限 & 构造跨空间、跨应用、越权读写、越权检索与越权问答场景，验证RBAC与content\_auth裁决结果。 & 越权内容不会出现在列表、检索与问答结果中；对象级与作用域级权限口径保持一致。 \\ \hline
    可扩展性 & 验证统一内容模型对knowledge、faq、course、case四类内容的承载能力，并检查治理规则与洞察指标扩展时对主链路的影响。 & 多形态内容共享统一元数据与生命周期管理；新增治理规则与指标维度主要落在扩展配置与子模块，未破坏主链路结构。 \\ \hline
    性能与稳定性 & 对高频读接口、发布后索引链路与问答链路进行分级压测，并观察错误率、消息积压、超时退化与链路恢复情况。 & 读链路延迟随QPS增长平稳上升；索引构建通过消息队列实现削峰；问答链路在大模型超时场景下可退化为检索摘要返回。 \\ \hline
  \end{tabularx}
\end{table}

\subsection{核心接口分级压测}

为验证平台在高频浏览、检索与问答场景下的可用性，本文对内容列表、内容检索、内容详情与智能问答接口进行分级压测。测试采用固定发布数据集与预热后的索引环境，逐步提升并发请求强度，并统一观察平均响应时间、P95、P99、错误率等指标。测试结果如表\ref{tab:perf_scope}所示。表中结果显示，内容列表、内容检索和内容详情接口在1000 QPS下P99分别为136ms、268ms和104ms，错误率均低于0.3\%；智能问答接口在30 QPS下首包P99为5110ms，主要耗时来自检索增强和大模型生成。

\begin{table}[!htbp]
  \centering
  \caption{核心接口分级压测结果表}
  \label{tab:perf_scope}
  \small
  \setlength{\tabcolsep}{4pt}
  \begin{tabularx}{\linewidth}{>{\raggedright\arraybackslash}p{0.14\linewidth}|c|c|c|c|>{\centering\arraybackslash}p{0.11\linewidth}|X}
    \hline
    接口 & QPS & 平均时延/ms & P95/ms & P99/ms & 错误率 & 备注 \\ \hline
    内容列表 & 500 & 31 & 61 & 89 & 0.04\% & 典型浏览负载 \\ \hline
    内容列表 & 1000 & 48 & 93 & 136 & 0.11\% &  \\ \hline
    内容检索 & 500 & 63 & 121 & 176 & 0.13\% & 含关键词与过滤条件 \\ \hline
    内容检索 & 1000 & 97 & 184 & 268 & 0.27\% &  \\ \hline
    内容详情 & 500 & 24 & 49 & 72 & 0.03\% & 详情页读取负载 \\ \hline
    内容详情 & 1000 & 36 & 71 & 104 & 0.06\% &  \\ \hline
    智能问答 & 15 & 1980 & 3010 & 3890 & 0.42\% & 流式输出的首包时延 \\ \hline
    智能问答 & 30 & 2720 & 3960 & 5110 & 0.89\% & 流式输出的首包时延 \\ \hline
  \end{tabularx}
\end{table}


\subsection{索引构建链路时延}

内容发布后，平台通过消息队列异步触发索引构建任务，分别写入全文索引和向量索引。本文以内容发布事件写入消息队列作为起点，以全文索引和向量索引均可被检索命中作为终点，统计索引构建端到端时延。测试结果如表\ref{tab:index_latency}所示。结果表明，索引构建链路在单条内容和批量发布场景下均能在分钟级内完成更新，适合平台以已发布内容为口径的检索与问答应用。

\begin{table}[!htbp]
  \centering
  \caption{索引构建链路端到端时延表}
  \label{tab:index_latency}
  \small
  \begin{tabularx}{\linewidth}{>{\raggedright\arraybackslash}p{0.20\linewidth}|c|c|c|c|X}
    \hline
    发布场景 & 内容数 & 平均时延/s & P50/s & P95/s & 备注 \\ \hline
    单条内容发布 & 1 & 3.2 & 2.7 & 6.8 & 包含切片与双索引写入 \\ \hline
    小批量发布 & 100 & 38.6 & 34.1 & 72.4 & \\ \hline
    批量导入发布 & 1000 & 286.5 & 251.8 & 498.6 & 按批次限流写入 \\ \hline
  \end{tabularx}
\end{table}

\subsection{RAG链路工程指标}

由于智能问答链路包含问题改写、检索召回、证据重排、大模型生成和引用映射等步骤，本文进一步从Token消耗、调用成本和引用正确率等工程指标进行评估。测试问题来自内容应用模块中的高频检索词和历史问答会话抽样，人工检查回答是否引用了与问题相关的已发布内容切片。相关结果如表\ref{tab:rag_metric}所示。结果表明，RAG链路能够在可控的Token消耗下返回带引用的回答，引用正确率达到工程应用的基本要求。

\begin{table}[!htbp]
  \centering
  \caption{RAG链路工程指标结果表}
  \label{tab:rag_metric}
  \small
  \begin{tabularx}{\linewidth}{>{\raggedright\arraybackslash}p{0.20\linewidth}|c|c|c|X}
    \hline
    指标 & 结果 & 样本规模 & 统计口径 & 备注 \\ \hline
    平均输入Token & 1850 & 300 & 证据集合与Prompt合计 &  \\ \hline
    平均输出Token & 326 & 300 & 首次完整回答 &  \\ \hline
    引用正确率 & 92.7\% & 300 & 人工抽检 & 引用内容需与回答结论直接相关 \\ \hline
  \end{tabularx}
\end{table}


\subsection{内容规模增长下的资源消耗}

为评估系统在内容规模增长时的扩展性，本文进一步从已发布内容规模为10万级与100万级两类场景出发，观察全文索引、向量索引以及服务侧关键资源的变化趋势。相关结果表如表\ref{tab:resource_scale}所示，结果表明，当内容规模从10万级增长到100万级时，ES占用从9.8GB增长到84.2GB，向量索引占用从15.6GB增长到129.4GB，检索P95从168ms增长到253ms，整体资源增长与内容规模增长基本匹配，检索链路仍处于可接受范围内。

\begin{table}[!htbp]
  \centering
  \caption{内容规模增长下的资源消耗结果表}
  \label{tab:resource_scale}
  \small
  \begin{tabularx}{\linewidth}{>{\raggedright\arraybackslash}p{0.14\linewidth}|>{\centering\arraybackslash}p{0.13\linewidth}|>{\centering\arraybackslash}p{0.15\linewidth}|>{\centering\arraybackslash}p{0.13\linewidth}|>{\centering\arraybackslash}p{0.12\linewidth}|X}
    \hline
    规模场景 & ES占用/GB & 向量索引占用/GB & 服务峰值内存/GB & 检索P95/ms & 备注 \\ \hline
    10万级内容 & 9.8 & 15.6 & 7.8 & 168 & 初始规模下的索引与检索链路表现 \\ \hline
    100万级内容 & 84.2 & 129.4 & 27.3 & 253 & 扩容后的资源变化与检索性能表现 \\ \hline
  \end{tabularx}
\end{table}


\subsection{应用效果评估}

除功能与性能指标外，平台上线后的价值还体现在内容运营效率和知识获取效率的改善。本文选取内容定位耗时、内容发布周期、治理任务处理周期和智能问答直接解决率等指标进行对比，评估对象为平台改造前后的同类内容管理与使用场景。相关结果如表\ref{tab:business_effect}所示。结果表明，平台通过统一内容组织、检索问答和治理任务线上化，降低了人工查找与协同处理成本，并提高了内容应用链路的响应效率。

\begin{table}[!htbp]
  \centering
  \caption{应用效果评估结果表}
  \label{tab:business_effect}
  \small
  \begin{tabularx}{\linewidth}{>{\raggedright\arraybackslash}p{0.24\linewidth}|c|c|c|X}
    \hline
    指标 & 改造前 & 改造后 & 变化 & 说明 \\ \hline
    内容定位平均耗时 & 6.8min & 2.1min & 降低69.1\% & 通过统一检索与类目标签定位内容 \\ \hline
    常规内容发布周期 & 2.5天 & 1.1天 & 缩短56.0\% & 审批流与发布口径线上化 \\ \hline
    治理任务处理周期 & 5.2天 & 2.4天 & 缩短53.8\% & 治理任务明确问题内容与责任人 \\ \hline
    高频问题直接解决率 & 61.5\% & 78.6\% & 提升17.1个百分点 & 由检索与RAG问答共同支撑 \\ \hline
  \end{tabularx}
\end{table}


\section{本章小结}

本章对平台的测试环境、单元测试、功能测试以及非功能与性能测试进行了说明。测试结果表明：平台在组织与权限边界、内容发布口径、事件驱动索引更新、RAG问答引用可追溯以及治理任务处理等关键链路上能够满足预期；在非功能维度上，平台满足一致性、安全隔离、可扩展性与稳定性要求；在性能维度上，高频读接口在1000 QPS下保持较低错误率，智能问答链路在30 QPS下能够稳定返回带引用结果；在应用效果方面，平台对内容定位、发布维护和治理处理效率均有改善。
